# 4. Minimal Theoretical Model

The purpose of this section is to construct the simplest dynamical model capable of representing nuclear-state control of fission.

The model is intentionally abstract. It does not assume the existence of a specific experimentally accessible pair of states in a particular fissile isotope. Instead, it identifies the minimum physical ingredients required for the proposed mechanism to operate and provides a quantitative framework that can later be instantiated using realistic nuclear data.

The minimal system contains two nuclear configurations,

[
|A\rangle
\qquad\text{and}\qquad
|B\rangle ,
]

which differ in their probability of undergoing fission.

We assume

[
\Gamma_f^{(A)}
\neq
\Gamma_f^{(B)},
]

where (\Gamma_f^{(i)}) denotes the fission decay rate associated with state (i).

An external actuator couples the two states,

[
|A\rangle
\overset{u_Q(t)}{\longleftrightarrow}
|B\rangle .
]

The central question is whether controlled redistribution of nuclear population between these states produces a measurable modification of the total fission probability.

---

## 4.1 Population-based two-state model

We first consider the simplest case in which quantum coherence between the two states is neglected.

Let

[
p_A(t)
]

and

[
p_B(t)
]

denote the populations of states (A) and (B).

The external actuator produces transitions with rates

[
W_{A\rightarrow B}(t)
]

and

[
W_{B\rightarrow A}(t).
]

Both states may decay through fission and through competing non-fission channels.

We therefore introduce

[
\Gamma_f^{(A)},\qquad
\Gamma_f^{(B)}
]

for fission and

[
\Gamma_d^{(A)},\qquad
\Gamma_d^{(B)}
]

for all competing decay processes.

The population dynamics is then

[
\frac{dp_A}{dt}
===============

*

W_{A\rightarrow B}p_A
+
W_{B\rightarrow A}p_B
---------------------

\Gamma_A p_A ,
]

[
\frac{dp_B}{dt}
===============

## W_{A\rightarrow B}p_A

## W_{B\rightarrow A}p_B

\Gamma_B p_B ,
]

where

[
\Gamma_A
========

\Gamma_f^{(A)}
+
\Gamma_d^{(A)}
]

and

[
\Gamma_B
========

\Gamma_f^{(B)}
+
\Gamma_d^{(B)}.
]

The instantaneous fission rate for a single initially prepared nucleus is

[
r_f(t)
======

\Gamma_f^{(A)}p_A(t)
+
\Gamma_f^{(B)}p_B(t).
]

The integrated probability that the system eventually undergoes fission is

[
P_f
===

\int_0^\infty
r_f(t),dt.
]

Without external control,

[
W_{A\rightarrow B}
==================

# W_{B\rightarrow A}

0,
]

and for an initially populated state (A),

[
p_A(0)=1,
\qquad
p_B(0)=0,
]

the fission probability is simply

[
P_f^{(0)}
=========

\frac{\Gamma_f^{(A)}}{\Gamma_A}.
]

The external actuator becomes physically relevant if it transfers population toward a state having a different fission branching ratio.

For example, if

[
\frac{\Gamma_f^{(B)}}{\Gamma_B}

>

\frac{\Gamma_f^{(A)}}{\Gamma_A},
]

then population transfer

[
A\rightarrow B
]

increases the integrated fission probability.

Conversely, if

[
\frac{\Gamma_f^{(B)}}{\Gamma_B}
<
\frac{\Gamma_f^{(A)}}{\Gamma_A},
]

the same control field suppresses fission.

This immediately establishes a minimal control criterion:

[
\boxed{
\frac{\Gamma_f^{(A)}}{\Gamma_A}
\neq
\frac{\Gamma_f^{(B)}}{\Gamma_B}
}
]

together with

[
\boxed{
W_{A\rightarrow B}
\neq0.
}
]

No long-lived quantum coherence is required.

---

## 4.2 Control efficiency

We define a dimensionless microscopic control efficiency

[
\eta_Q
======

\frac{
P_f^{(Q)}-P_f^{(0)}
}{
P_f^{(0)}
},
]

where (P_f^{(Q)}) denotes the fission probability in the presence of the external actuator.

Thus,

[
\eta_Q>0
]

corresponds to enhancement of fission,

while

[
\eta_Q<0
]

corresponds to suppression.

The magnitude

[
|\eta_Q|
]

provides a natural first measure of the effectiveness of nuclear-state control.

A value such as

[
|\eta_Q|\ll1
]

would indicate that the microscopic effect is probably irrelevant at reactor scale unless extraordinarily large numbers of nuclei can be controlled.

Conversely, an effect approaching

[
|\eta_Q|
\sim
10^{-2}-1
]

would immediately motivate investigation of macroscopic consequences.

The purpose of introducing (\eta_Q) is not to assume that such large values can be obtained, but to define a quantity that can be calculated from microscopic nuclear models and ultimately measured experimentally.

---

## 4.3 Competition between control and relaxation

Population transfer must occur before the nuclear system loses memory of the externally prepared state.

We therefore introduce a characteristic control timescale

[
\tau_Q
\sim
\frac{1}{W_Q}
]

and a characteristic intrinsic relaxation or equilibration timescale

[
\tau_{\mathrm{eq}}.
]

For nuclear-state manipulation to influence fission before equilibration erases the prepared state, a necessary condition is approximately

[
\boxed{
\tau_Q
\lesssim
\tau_{\mathrm{eq}}
}
]

or equivalently

[
\boxed{
W_Q
\gtrsim
\Gamma_{\mathrm{eq}}.
}
]

This relation may prove to be one of the most restrictive conditions in realistic fissile nuclei.

If instead

[
\tau_Q
\gg
\tau_{\mathrm{eq}},
]

the externally induced redistribution is erased before it can affect the subsequent fission dynamics.

The competition

[
\text{external control}
\quad\text{vs.}\quad
\text{intrinsic equilibration}
]

is therefore a central dimensionless parameter of the theory.

We define

[
\mathcal C_Q
============

\frac{W_Q}{\Gamma_{\mathrm{eq}}}.
]

The regimes

[
\mathcal C_Q\ll1,
]

[
\mathcal C_Q\sim1,
]

and

[
\mathcal C_Q\gg1
]

correspond respectively to ineffective, competing and dominant external control.

---

## 4.4 Coherent two-state model

The population model can be extended to include genuine quantum coherence.

Consider the Hamiltonian

[
H(t)
====

\frac{\hbar}{2}
\begin{pmatrix}
-\Delta & \Omega(t)e^{-i\phi(t)}
\
\Omega(t)e^{i\phi(t)} & \Delta
\end{pmatrix},
]

where

[
\Omega(t)
]

is the externally controlled Rabi coupling,

[
\Delta
]

is the detuning between the applied field and the nuclear transition,

and

[
\phi(t)
]

is the controllable phase.

The nuclear state is described by a density matrix

[
\rho
====

\begin{pmatrix}
\rho_{AA} & \rho_{AB}
\
\rho_{BA} & \rho_{BB}
\end{pmatrix}.
]

Its evolution may be represented schematically by a master equation

[
\frac{d\rho}{dt}
================

-\frac{i}{\hbar}[H,\rho]
+
\mathcal L_{\mathrm{decay}}[\rho]
+
\mathcal L_{\mathrm{decoh}}[\rho].
]

The first term describes coherent external control.

The second represents irreversible nuclear decay, including fission and competing channels.

The third represents decoherence and intrinsic equilibration.

The instantaneous fission rate becomes

[
r_f(t)
======

\Gamma_f^{(A)}\rho_{AA}(t)
+
\Gamma_f^{(B)}\rho_{BB}(t)
+
r_f^{(\mathrm{int})}(t),
]

where

[
r_f^{(\mathrm{int})}
]

represents any interference contribution between amplitudes leading from (A) and (B) into common fission channels.

If such common channels exist, one may schematically write

[
r_f^{(\mathrm{int})}
\propto
2,
\mathrm{Re}
\left[
\rho_{AB}
A_{A\rightarrow f}
A_{B\rightarrow f}^{*}
\right].
]

The important consequence is that the fission probability may then depend not only on populations but also on phase.

Thus,

[
P_f
===

P_f(\Omega,\Delta,\phi).
]

In such a regime, changing the phase of an external field could modify the fission probability even without changing the total excitation deposited in the nucleus.

Observation of

[
\frac{\partial P_f}{\partial\phi}\neq0
]

would therefore constitute a particularly strong signature of genuine coherent control of fission.

---

## 4.5 Three physically distinct regimes

The minimal model naturally identifies three regimes.

### Regime I — No nuclear-state control

If

[
W_Q
\ll
\Gamma_{\mathrm{eq}},
]

the externally driven dynamics is negligible.

The nuclear-state distribution rapidly approaches its natural statistical form, and

[
P_f^{(Q)}
\simeq
P_f^{(0)}.
]

This corresponds to conventional reactor physics.

### Regime II — Population control

If external excitation produces significant redistribution among nuclear states before decay or equilibration,

[
W_Q
\gtrsim
\Gamma_{\mathrm{eq}},
]

but coherence is rapidly destroyed, the density matrix becomes approximately diagonal.

Then

[
P_f
===

\sum_i p_i P_f^{(i)}.
]

Control exists, but it is effectively classical at the ensemble level.

This is the weakest and potentially most experimentally accessible form of nuclear-state fission control.

### Regime III — Coherent control

If

[
\Omega
\gtrsim
\Gamma_{\mathrm{decoh}}
]

for sufficiently long times, off-diagonal density-matrix elements survive,

[
\rho_{AB}\neq0.
]

The fission probability may then depend on quantum interference and the phase of the external actuator.

This regime constitutes the closest nuclear analogue to coherent control techniques used in atomic and molecular physics.

It is also the most demanding physical regime.

---

## 4.6 Coupling the microscopic model to a neutron field

The previous model describes the response of a single nuclear ensemble.

To connect it to reactor physics, consider a neutron flux

[
\Phi(E,t).
]

Instead of assigning a fixed microscopic fission cross section to the fuel, we define

[
\sigma_f^{\mathrm{eff}}
=======================

\sigma_f^{\mathrm{eff}}
(E,\rho_N).
]

For the two-state population model,

[
\sigma_f^{\mathrm{eff}}(E,t)
============================

p_A(t)\sigma_f^{(A)}(E)
+
p_B(t)\sigma_f^{(B)}(E).
]

The macroscopic fission rate becomes

[
R_f(t)
======

N_f
\int
\Phi(E,t)
\left[
p_A(t)\sigma_f^{(A)}(E)
+
p_B(t)\sigma_f^{(B)}(E)
\right]
dE.
]

The external control therefore enters the reactor equations through

[
p_i=p_i[u_Q(t)].
]

This can be represented schematically as

[
u_Q(t)
\rightarrow
p_i(t)
\rightarrow
\sigma_f^{\mathrm{eff}}(t)
\rightarrow
R_f(t).
]

The conventional and quantum-control channels coexist:

[
\Phi(E,t)
]

remains a dynamical variable determined by neutron transport, while

[
\sigma_f^{\mathrm{eff}}(E,t)
]

acquires an additional controlled dependence on the nuclear state.

---

## 4.7 Effective reactivity contribution

To connect the microscopic effect directly to reactor kinetics, suppose that external control induces

[
\sigma_f
\rightarrow
\sigma_f+\delta\sigma_f.
]

The corresponding perturbation of the macroscopic fission cross section is

[
\delta\Sigma_f
==============

N_f\delta\sigma_f.
]

For a small perturbation, the resulting change in the effective multiplication factor may be written formally as

[
\delta k_{\mathrm{eff}}
\simeq
\frac{\partial k_{\mathrm{eff}}}
{\partial\Sigma_f}
\delta\Sigma_f,
]

with all other system variables held fixed.

The associated reactivity perturbation is

[
\delta\rho_Q
\simeq
\frac{\delta k_{\mathrm{eff}}}
{k_{\mathrm{eff}}^2}.
]

At criticality,

[
k_{\mathrm{eff}}\simeq1,
]

so approximately

[
\delta\rho_Q
\simeq
\delta k_{\mathrm{eff}}.
]

The quantity

[
\boxed{
\rho_Q
}
]

can therefore be interpreted as an effective reactivity contribution arising from nuclear-state manipulation.

This does not imply that quantum reactivity is additive under arbitrary conditions. The definition is intended as a first-order perturbative quantity suitable for comparison with conventional reactivity effects.

---

## 4.8 Extension of point-reactor kinetics

At the macroscopic level, ordinary point kinetics may be written as

[
\frac{dn}{dt}
=============

\frac{\rho-\beta}{\Lambda}n
+
\sum_i \lambda_i C_i,
]

where (n) denotes the neutron population.

In the present framework,

[
\rho
====

\rho_{\mathrm{conv}}
+
\rho_Q.
]

The neutron equation therefore becomes

[
\frac{dn}{dt}
=============

\frac{
\rho_{\mathrm{conv}}
+
\rho_Q
------

\beta
}{
\Lambda
}
n
+
\sum_i \lambda_i C_i.
]

However, unlike a conventional externally prescribed reactivity insertion,

[
\rho_Q
]

has its own microscopic dynamics.

Schematically,

[
\rho_Q
======

\rho_Q[\rho_N(t)],
]

while

[
\frac{d\rho_N}{dt}
==================

\mathcal F
[
\rho_N,
u_Q,
\Phi,
T,\ldots
].
]

The combined system therefore takes the form

[
\boxed{
\begin{aligned}
\dot{\rho}*N
&=
\mathcal F(\rho_N,u_Q,\Phi,\ldots),
[4pt]
\rho_Q
&=
\mathcal G(\rho_N),
[4pt]
\dot n
&=
\frac{
\rho*{\mathrm{conv}}+\rho_Q-\beta
}{\Lambda}n
+
\sum_i\lambda_iC_i.
\end{aligned}
}
]

This represents the minimal mathematical architecture of a reactor whose nuclear fuel possesses an externally controllable internal state.

---

## 4.9 Feedback between neutron field and nuclear state

The model also reveals an important point that is absent from a simple external-reactivity picture.

The coupling may be bidirectional.

The external field modifies the nuclear state,

[
u_Q
\rightarrow
\rho_N,
]

but the neutron field itself may also modify that state through excitation, absorption and compound-nucleus formation:

[
\Phi
\rightarrow
\rho_N.
]

At the same time,

[
\rho_N
\rightarrow
\sigma_f^{\mathrm{eff}}
\rightarrow
\Phi.
]

The complete structure is therefore a feedback loop:

[
\boxed{
u_Q
\rightarrow
\rho_N
\rightarrow
\sigma_f^{\mathrm{eff}}
\rightarrow
\Phi
\rightarrow
\rho_N
}
]

rather than a simple one-way actuator.

This feedback may generate dynamical behavior that has no exact analogue in conventional reactor control.

Whether such behavior is relevant depends entirely on the strength and timescale of the microscopic coupling.

---

## 4.10 Minimal falsifiable prediction

The proposed framework admits a simple falsification criterion.

Consider two otherwise identical experiments on a fissile target exposed to the same incident neutron distribution,

[
\Phi_1(E)
=========

\Phi_2(E).
]

In experiment 1, no state-control field is applied.

In experiment 2, an external control protocol prepares a different nuclear-state distribution while leaving the incident neutron field unchanged.

The hypothesis predicts that, under suitable conditions,

[
\rho_N^{(1)}
\neq
\rho_N^{(2)}
]

should produce

[
P_f^{(1)}
\neq
P_f^{(2)}.
]

Equivalently,

[
\boxed{
\left.
\frac{\partial P_f}
{\partial u_Q}
\right|_{\Phi}
\neq0.
}
]

This derivative provides perhaps the simplest mathematical statement of the entire proposal.

The question is not merely whether radiation can itself induce fission.

Rather, one asks whether an external interaction can modify the subsequent fission response to an otherwise unchanged neutron field.

A particularly strong experimental signature would be obtained if the modification depended on a control parameter such as pulse timing, detuning or phase,

[
P_f
===

P_f(u_Q),
]

while the total neutron exposure and deposited excitation energy were held fixed.

---

## 4.11 What the minimal model establishes

The model does not demonstrate that nuclear-state control is technologically feasible.

It establishes instead that the concept can be expressed as a well-defined dynamical problem.

The minimum ingredients are:

1. at least two nuclear configurations with distinguishable fission responses;

2. an externally controllable transition or redistribution mechanism between them;

3. a control rate competitive with intrinsic relaxation;

4. a non-zero change in the ensemble-averaged fission probability;

5. for reactor relevance, a sufficiently large induced modification of the macroscopic fission cross section.

The hypothesis therefore reduces to a quantitative search for the parameter domain

[
\boxed{
\left{
\Delta P_f,,
W_Q,,
\Gamma_{\mathrm{eq}},,
f_Q
\right},
}
]

where (f_Q) denotes the fraction of fissile nuclei effectively addressed by the external actuator.

If no physically accessible nuclear system satisfies these conditions, nuclear-state control cannot become relevant to reactor physics.

If such a system exists, however, the microscopic nuclear state becomes an additional dynamical variable in the theory of controlled fission.
